% !TeX root = ../papers/nc-p3.tex
\section{中心指標の倍率安定群、指数の付値と多項式次数の正則判定}
\label{sec:central-stabilizer-strata}
中心有限な標準正則折り畳みを $B$、頂点隅を $R$、中心を $Z$、一般の頂点可除環の
次数を $s$ とする。中心では経路長4を次数1と数える。
一位置の矢だけを倍率変更する $\mathbb G_m$ はこの次数付けを誘導する。

\begin{theorem}\label{thm:central-stabilizer-strata}
中心上の倍率作用は忠実で、非原点の幾何学的閉点の安定群の位数 $m$ は $s$ を割る。
四辺有限で $s>1$ なら $m<s$ である。
任意の $d\ge1$ について頂点隅の $d$ 次Veroneseの中心は $Z^{(d)}$ で、
その一般可除指数は同じ $s$ である。
\end{theorem}
\begin{proof}
$B$ の両側Koszul分解の自己双対性は内部移動4、ホモロジー移動4を持つ
（補題\ref{lem:cy-finite-birational-transfer}の証明）。
従って $HH^0(B)$ の単位に対応する非零の $HH_4(B)$ の類は経路長4を持つ。
最終Koszul項の各成分は四つの矢位置をちょうど一度通るテンソル直線と
次数零の頂点冪等元からなる最小次数部分を持つ。よってそのサイクルは倍率
$\lambda$ により $\lambda$ 倍される。$Q=M_4(D)$ への中心局所化後も消えない。
実際、両端を局所化した分解は $Q\otimes_B Q=Q$ の両側射影分解であり、
Hochschild複体の最終項では各非零係数の冪等元が残り、流入する境界もない。
$Z$ を各点ごとに固定する倍率は $Q$ 上内部自己同型なので $HH_4(Q/k)$ に恒等に作用し、
$\lambda=1$ となる。従って一般安定群は自明である。さらに
$K=\operatorname{Frac}Z$ 上の $D$ に $G=\mu_d$ が作用し、$K/K^G$ は群 $G$ のGalois拡大。
Galois降下から $D^G$ は次数 $s$ の中心的単純環で、零因子を持たない。
中心分母の全共役の積により $(Z^G\setminus0)^{-1}R^G=D^G$ となるので
$Z(R^G)=R^G\cap K^G=Z^G$ である。ここで $R\cap K=Z$ は
$R$ の被約traceが $Z$ に入り、中心分数上の元のtraceは $s$ 倍となることによる。

非原点 $x$ では非零な正次数値があるので安定群は有限巡回群 $\mu_m$ である。
被約traceの半単純特殊化の各因子は正重複度を持ち、$\mu_m$ は重複度を保って因子を置換する。
一因子の安定群を $\mu_d$ とし、その生成元との絡作用素を $d$ 乗が1となるよう取る。
固有空間は、倍率を変えた矢位置を通るごとに一つ進み、$4d$ 個の頂点を巡る表示を与える。
$x$ で非零な斉次中心元 $z$ の次数は $m$ の倍数である。各始点成分を
$4d$ 周期の各持ち上げで同じ経路の式として取れば、終点も同じ頂点になり中心元を得る。
中心性は各矢との可換式を経路の始点を指定して持ち上げればよい。
この元は持ち上げた表示に非零スカラーとして作用し、各頂点単純には零として作用する。
よってそれらの間の全Extは消える。各頂点のKoszul分解のEuler式により
固有空間次元列の四階差は零となる。周期的なので全項は同じ整数 $c\ge1$ である。
従って元の各頂点の次元は $dc$、軌道長は $m/d$ であり、軌道の寄与は
$m c$ の重複度倍である。総和は $s$ だから $m\mid s$ を得る。

四辺有限で $m=s>1$ なら一軌道かつ $c=1$、重複度1である。
$d>1$ なら四移動の恒等性により一次元固有空間から組み直した元の単純は
薄い真部分表示を持つ。$d=1$ なら同一グラフ点の異なる倍率の薄い単純だけで
繊維を尽くし、第一Extの消滅と繊維の連結性に反する。いずれも
定理\ref{thm:odd-prime-finite-indices}の最後の議論で、ここには素数性を使わない。
\end{proof}

\begin{corollary}\label{cor:central-valuation-strata}
素数 $p$ について $r_j=\dim(\operatorname{Spec} Z)^{\mu_{p^j}}$ と置くと
\[
 2v_p(s)\le v_p(64)+\sum_{j\ge1}r_j .
\]
四辺有限の場合、奇素数の二乗を一般可除指数に持つ正則窓は存在しない。
同じ仮定で指数が奇素数 $p$ の三乗なら $r_1=r_2=3$ が必要である。
\end{corollary}
\begin{proof}
$r_j$ は降列で、忠実性から $r_1\le3$、$j>v_p(s)$ なら原点だけなので $r_j=0$。
$k=1,\ldots,4$ について
$b_k=\min\{b\ge0:r_{b+1}\le4-k\}$ とする。$b=b_k$ として $Z$ の有限斉次生成元のうち
正次数の付値が $b$ 以下のものの共通零点はちょうど $\mu_{p^{b+1}}$ の固定集合である。
これらの冪を共通次数 $n_k$ に揃え、さらに $p$ 冪を掛けて $v_p(n_k)=b$ とできる。
逆に次数 $n_k$ の元は全てその固定集合に消える。従って一般のその次数の元は、先行する
$k-1$ 個の切断で残った次元 $5-k$ のどの成分も含まない。順に四つ選べば
多項式パラメータ環を得る。$\dim R_n=\binom{4n+3}3$ からその上の階数は
$64\prod n_k$ であり $s^2$ の倍数である。$b_k\ge j$ は $r_j>4-k$ と同値なので
$\sum_k b_k=\sum_j r_j$ となる。四辺有限かつ $s=p^a$、$a\ge1$ なら $r_a=0$。
奇素数で $a=2$ なら $4\le3$ が必要となり、$a=3$ なら $6\le r_1+r_2\le6$ となる。
\end{proof}

\begin{corollary}\label{cor:polynomial-pi-window-cutoff}
定理\ref{thm:noetherian-pi-window-cutoff}の全有限Frobenius候補に対する同値判定では、
一般頂点可除指数の上界を $u$ として、$1\le s\le u$ のいずれかに対する経路長
$h=16s^4$ の中心四元の検査と、有限乗法表での指数上界の検査を使える。
検査の経路次数は高々 $64u^4$ である。
\end{corollary}
\begin{proof}
必要性では $s$ を実際の一般指数とする。非原点中心指標 $x$ で $R$ の被約traceの
半単純特殊化を取る。その像 $E\subset M_s(k)$ は半単純、通常のtrace対合は
各単純因子上で正の整数重複度を掛けた行列trace対合なので非退化である。
次数1の生成元の語で単位から作る累積張空間は安定すると $E$ だから、長さ
$s^2-1$ までの語に基底がある。$D_0=2s^2-1$ とする。倍率 $\lambda$ による
$x$ の移動後に長さ $D_0$ までの語のtraceが全て同じと仮定する。像の次元は移動前後で同じであり、
基底語のGram行列は等しく可逆である。基底語二つの間に一生成元を入れた
traceも等しいので、全生成元の左乗法の係数が一致する。単位を基底に含めれば
全語のtraceが一致する。中心元のtraceはその値の $s$ 倍なので $\lambda x=x$ である。

よって $x$ で非零値を持つ長さ $1$ から $D_0$ の語の被約traceの次数集合は空でなく、
その最大公約数はちょうど安定群の位数 $m$ である。被約trace自体は $Z$ の斉次元である。
正整数の集合を $m$ で割ると、最小元 $a$ を法とする全剰余類を生成する。
各剰余類は、加法の剰余グラフの単純路により高々 $a-1$ 個の生成元の和で表せる。
そこに $a$ を加えていけば、元の次数で少なくとも $D_0^2$ の全ての $m$ の倍数を表せる。
$m\mid s$ なので $4s^4$ はこの範囲にある。従ってその次数の中心元は全非原点で
共通には消えない。一般四元で切る。これらは元の環ですでに中心なので、多項式環を使う定理\ref{thm:onesided-parameter-cutoff}を適用できる。
その $D=4h$ までの経路次元とKoszul複体、および $D-3$ までの剰余で足り、逆も同じ定理による。
\end{proof}
